% ABSTRACT page — Juknis p.3: heading, title repeated (italic),
% "By / Name / NIM / Program", then body (italic, not bold, 12pt),
% keywords (max. 7) separated by two blank lines from the body.
\chapter*{\textit{Abstract}}
\addcontentsline{toc}{chapter}{Abstract}

% Metadata from config.tex — edit that file, not this one.
\begin{center}
{\bfseries\itshape\fontsize{14}{14}\selectfont \ThesisTitleEN}

\vspace{1\baselineskip}
\textit{By}\\
\textbf{\ThesisAuthor}\\
NIM: \ThesisNIM\\
\textit{(\ThesisProgramEN)}
\end{center}

\vspace{1\baselineskip}

\itshape
Facial Expression Recognition (FER) models are usually developed on benchmark data that do not represent all groups, so their performance can differ between groups of users. This thesis evaluates such a performance gap between two skin-brightness groups, Dark and Medium-Dark, on 135 manually labeled face crops (five classes) from one Papuan documentary film. Brightness groups are computed from pixel luminance and treated as a proxy, not as a validated measure of skin tone or ethnicity. Three research questions are answered by re-analyzing saved predictions, with 95\% confidence intervals (CI) from a grouped bootstrap. First, the pretrained HSEmotion model reaches an accuracy of 0.363 and a macro F1 of 0.249; its accuracy is 0.483 on Dark crops and 0.130 on Medium-Dark crops, a gap of 0.353 (95\% CI [0.182, 0.540]). The gap persists within the same class and comes largely from Neutral expressions predicted as Angry in the Medium-Dark group (12 of 21). Second, one exploratory adaptive CLAHE configuration changes the gap by $-$0.023 (95\% CI [$-$0.136, 0.073]); its effect cannot be determined from this sample. Third, in a retrospective comparison, a Logistic Regression classifier on frozen ArcFace embeddings reaches an accuracy of 0.407 and a macro F1 of 0.272, with no significant difference in overall accuracy from HSEmotion (McNemar's test, $p = 0.4614$). Its gap is smaller, 0.123 (95\% CI [$-$0.038, 0.282]), because of higher Medium-Dark accuracy; however, the difference between the two models' gaps excludes zero under only one of the two groupings examined, so this reduction is suggestive, not established. These results are associations in one small sample from one film with one annotator. This thesis does not evaluate lighter skin, does not infer causation, and does not generalize its results to other populations. A facial-landmark annotation track (Papuan V2) is documented as supporting work and is not used as evidence.

\vspace{2\baselineskip}
\upshape
Keywords: facial expression recognition, fairness evaluation, group performance gap, skin brightness (proxy), HSEmotion, CLAHE, ArcFace.
